\cvsection{Employment and Research}

\vspace{-0.5mm}

\begin{cventries}

\cventry
{Graduate Research Assistant}
{University of Hawai`i at Mānoa}
{Honolulu, HI}
{August 2022 - present, full-time}
{
    Research \textbf{probabilistic machine learning} with Dr. Peter Sadowski, applying \textbf{Gaussian Processes} in \textbf{Python} for anomaly detection and accelerating \textbf{simulation-based inference} via \textbf{neural likelihood models} [\ref{neural_surrogate}].
}

\vspace{0mm}

\cventry
{SW-GPU Intern, PhysicsNeMo / Earth-2}
{NVIDIA}
{Santa Clara, CA}
{July 2026 - October 2026, full-time}
{
    Researched deep learning applications for climate science, developing generative satellite and radar weather models.
}

\vspace{0mm}

\cventry
{Research Intern, MLIA}
{NASA Jet Propulsion Laboratory, Caltech}
{Pasadena, CA}
{March 2026 - June 2026, full-time}
{
    Optimized ionospheric deep learning architectures to reduce runtime and built evaluation pipelines across geomagnetic storm conditions [\ref{ionospheric_forecasts}].
}

\vspace{0mm}

\cventry
{Researcher at HelioLab}
{Frontier Development Lab}
{Remote}
{June 2025 - October 2025, full-time}
{
    Developed a novel dataset and deep learning models in \textbf{PyTorch} to forecast ionospheric dynamics, presented at NeurIPS [\ref{ionosphere_dataset}, \ref{ioncast}].
}

\vspace{0mm}

\cventry
{Data Science Intern}
{Striveworks, Inc}
{Austin, TX}
{May 2023 - August 2023, full-time}
{
    Researched \textbf{large language models} for NLP tasks, evaluating parameter-efficient fine-tuning strategies such as LoRA.
}

\vspace{0mm}

\cventry
{Undergraduate Research Assistant}
{Laboratory for Atmospheric and Space Physics}
{Boulder, CO}
{August 2020 - July 2022, part-time}
{
    Implemented K-means and \textbf{CNNs} in \textbf{PyTorch} to segment solar polar coronal holes [\ref{search}].
}

\vspace{0mm}

\cventry
{Research Experience for Undergraduates}
{MIT Haystack Observatory}
{Westford, MA}
{May 2021 - August 2021, full-time}
{
    Applied GMMs and \textbf{CNNs} to detect ocean wave events in seismic data from Antarctica's Ross Ice Shelf.
}

\vspace{0mm}

\cventry
{Research Experience for Undergraduates}
{Institute for Astronomy at the University of Hawai`i at Mānoa}
{Honolulu, HI}
{May 2020 - August 2020, full-time}
{
    Analyzed the entire Kepler Space Telescope dataset using \textbf{Python} statistical tools such as \textbf{Pandas} to evaluate the biases in its selection function [\ref{kepler}].
}

\vspace{-2mm}

\end{cventries}
